\documentclass[11pt]{article}
\usepackage[margin=1in]{geometry}
\usepackage{amsmath,amssymb,mathtools}
\usepackage{fontspec,unicode-math}
\setmainfont{Latin Modern Roman}
\setsansfont{Latin Modern Sans}
\setmathfont{Latin Modern Math}
\usepackage{xurl,hyperref}
\hypersetup{colorlinks=true,urlcolor=blue,linkcolor=blue}
\newcommand{\sref}[2]{\hyperlink{source:#1:#2}{[S#2]}}
\newcommand{\cataloguescope}[1]{\paragraph{Recorded mathematical question.}#1}
\begin{document}
\section{Open Problems on Billiards and Geometric Optics}\label{statement}

\subsection{Definitions and mathematical statement}

\paragraph{Shared notation.}$\mathbb N=\{1,2,\ldots\}$, and $\mathbb Z,\mathbb Q,\mathbb R,\mathbb C$ denote the integers, rationals, reals and complex numbers. Any other conventions are specified locally.


An oval is a smooth closed strictly convex plane mirror. Fix a point source in its interior. For each integer $n\ge1$, the $n$th caustic by reflection is the envelope of the one-parameter family of rays after $n$ specular reflections. Cusps are the ordinary cusp singularities of a generic envelope; the source restricts to generic caustics and source positions in general position.
\paragraph{Statement recorded in the supplied source.}
For an oval $\gamma$ and a light source inside it, call the envelope of rays after $n$ reflections the $n$th caustic by reflection. Is every generic caustic by reflection in an ellipse a curve with exactly four cusps? Does this four-cusp property characterize ellipses? \sref{statement}{2}
\subsection{Short English statement}
Do all generic reflected caustics of every order in an ellipse have exactly four cusps, and does having this property for interior light sources characterize ellipses?
\subsection{Sources}
\begin{description}
\item[\hypertarget{source:statement:1}{S1}] ULAM.AI and the UnsolvedMath Contributors, \emph{UnsolvedMath: A Curated Collection of Open Mathematics Problems}, catalogue entry. CC BY 4.0. \url{https://huggingface.co/datasets/ulamai/UnsolvedMath}.
\item[\hypertarget{source:statement:2}{S2}] M. Bialy, C. Fierobe, A. Glutsyuk, M. Levi, A. Plakhov and S. Tabachnikov, \emph{Open problems on billiards and geometric optics} (2021), arXiv:2110.10750v2, Serge Tabachnikov, Problem 4, p. 12. \url{https://arxiv.org/pdf/2110.10750v2}.
\end{description}
\label{end:statement}
\end{document}
